\documentclass[a4paper]{arrowhead}

\usepackage[yyyymmdd]{datetime}
\usepackage{etoolbox}
\usepackage[utf8]{inputenc}
\usepackage[T1]{fontenc}
\usepackage{multirow}
\usepackage{longtable}
\usepackage{tikz}
\usetikzlibrary{positioning,shapes.multipart,arrows.meta}

\renewcommand{\dateseparator}{-}
\setlength{\emergencystretch}{3em}
\lstset{language=HTTP,breakatwhitespace=false,columns=fullflexible}

%% Special references
\newcommand{\mref}[2][]{{\textcolor{ArrowheadPurple}{\hyperref[sec:model:#2]{\ifstrempty{#1}{#2}{#1}}}}}
\newcommand{\pref}[1]{{\textcolor{ArrowheadGrey}{\hyperref[sec:model:primitives:#1]{#1}}}}
%%

\begin{document}
\sloppy

%% Arrowhead Document Properties
\ArrowheadTitle{AIGym}
\ArrowheadType{System Description}
\ArrowheadTypeShort{SysD}
\ArrowheadVersion{5.2.2}
\ArrowheadDate{2026-10-30}
\ArrowheadAuthor{Jerker Delsing}
\ArrowheadStatus{RELEASE CANDIDATE}
\ArrowheadContact{jerker.delsing@ltu.se}
\ArrowheadFooter{\href{www.arrowhead.eu}{www.arrowhead.eu}}
\ArrowheadSetup
%%

%% Front Page
\begin{center}
  \vspace*{1cm}
  \huge{\arrowtitle}

  \vspace*{0.2cm}
  \LARGE{\arrowtype}
  \vspace*{1cm}

  \vspace*{\fill}

  % Front Page Image

  \vspace*{1cm}
  \vspace*{\fill}

  % Front Page Abstract
  \begin{abstract}
    This document describes the AIGym support system of the Eclipse Arrowhead 5.2.2 framework, as implemented in \texttt{ah5-AIGym-java-spring} of the primary Java source tree (\texttt{Eclipse-Arrowhead-java}) and deployed with the standalone Compose definition of that folder.
  \end{abstract}

  \vspace*{1cm}
\end{center}

\newpage
%%

%% Table of Contents
\tableofcontents
\newpage
%%

\section{Overview}
\label{sec:overview}

This document describes the AIGym system. AIGym is an experimentation system for machine learning in the local cloud: it collects training samples, predicts with k nearest neighbours and consumes the \texttt{sysMonitorData} stream.

\begin{figure}[ht!]
  \centering
  \begin{tikzpicture}[font=\small, >=Stealth]
    \node[draw=ArrowheadBlue, very thick, minimum width=4.2cm, minimum height=2.2cm, align=center, fill=ArrowheadBlue!10] (s)
      {\guillemotleft{}System\guillemotright{}\\\textbf{AIGym}};
    \node[draw=ArrowheadTeal, rectangle, align=left, left=1.3cm of s, text width=4.6cm, inner sep=5pt] (p)
      {\textbf{Produced services}\\ \texttt{aiGym}\\ \texttt{sysMonitor}};
    \node[draw=ArrowheadOrange, rectangle, align=left, right=1.3cm of s, text width=4.6cm, inner sep=5pt] (c)
      {\textbf{Consumed}\\ AMQP broker (SysMonitor)\\ Authentication\\ ConsumerAuthorization\\ DynamicServiceOrchestration\\ ServiceRegistry\\ SysMonitor};
    \draw[->, thick, ArrowheadTeal] (s) -- (p);
    \draw[->, thick, ArrowheadOrange] (s) -- (c);
  \end{tikzpicture}
  \caption{SysML-style block view of AIGym with its produced and consumed services.}
  \label{fig:system}
\end{figure}

The rest of this document is organized as follows.
In Section \ref{sec:prior_art}, we reference major prior art capabilities of the system.
In Section \ref{sec:use}, we describe the intended usage of the system.
In Section \ref{sec:properties}, we describe fundamental properties provided by the system.
In Section \ref{sec:delimitations}, we describe delimitations of capabilities of the system.
In Section \ref{sec:services}, we describe the abstract service functions consumed or produced by the system.
In Section \ref{sec:security}, we describe the security capabilities of the system.

\newpage

\subsection{Significant Prior Art}
\label{sec:prior_art}

New in this workspace.

\subsection{How This System Is Meant to Be Used}
\label{sec:use}

Clients add samples and request predictions; AIGym orchestrates \texttt{sysMonitorData} and keeps the latest observations.

\subsection{System functionalities and properties}
\label{sec:properties}

\subsubsection{Functional properties of the system}

\begin{itemize}
  \item k-NN model store.
  \item sysMonitorData subscription.
  \item sysMonitor publication with \texttt{trainingSample} and \texttt{inference} observations.
\end{itemize}

\subsubsection{Configuration of system properties}

The system is configured with Spring properties: the bundled \texttt{application.\allowbreak{}yml} reads environment variables, and the standalone deployment overrides it with \texttt{config/\allowbreak{}application.\allowbreak{}properties} (copied into the image, \texttt{spring.\allowbreak{}config.\allowbreak{}additional-\allowbreak{}location}), whose passwords come from Compose file secrets. The properties that shape its behaviour are:

\noindent\begin{longtable}{| p{5.4cm} | p{3.4cm} | p{6.4cm} |} \hline
\rowcolor{gray!33} Property (variable) & Default & Meaning \\ \hline \endhead
\texttt{server.\allowbreak{}port} (\texttt{PORT}) & 8478 & HTTPS port (mutual TLS, \texttt{client-\allowbreak{}auth=\allowbreak{}need}). \\ \hline
\texttt{arrowhead.\allowbreak{}advertised-\allowbreak{}host} (\texttt{ADVERTISED\_\allowbreak{}HOST}) & aigym & Registered address. \\ \hline
\texttt{arrowhead.\allowbreak{}registry-\allowbreak{}url} (\texttt{SERVICE\_\allowbreak{}REGISTRY\_\allowbreak{}URL}) & \texttt{https:\allowbreak{}/\allowbreak{}/\allowbreak{}serviceregistry:\allowbreak{}8443} & The only configured Arrowhead endpoint; Authentication and the orchestrator are looked up there. \\ \hline
\texttt{arrowhead.\allowbreak{}authentication-\allowbreak{}url} (\texttt{AUTHENTICATION\_\allowbreak{}URL}) & (empty) & Optional override of the Authentication lookup. \\ \hline
\texttt{arrowhead.\allowbreak{}orchestration-\allowbreak{}ttl-\allowbreak{}seconds} (\texttt{ORCHESTRATION\_\allowbreak{}TTL\_\allowbreak{}SECONDS}) & 300 & Reuse of orchestration results. \\ \hline
\texttt{arrowhead.\allowbreak{}mqtt-\allowbreak{}host} / \texttt{-port} (\texttt{MQTT\_\allowbreak{}HOST} / \texttt{MQTT\_\allowbreak{}PORT}) & mqtt-broker / 8883 & sysMonitorData broker. \\ \hline
\texttt{arrowhead.\allowbreak{}amqp-\allowbreak{}host} / \texttt{-port} / \texttt{-exchange} (\texttt{AMQP\_\allowbreak{}*}) & amqp-broker / 5671 / arrowhead & sysMonitor publication. \\ \hline
\texttt{arrowhead.\allowbreak{}status-\allowbreak{}interval-\allowbreak{}ms} (\texttt{STATUS\_\allowbreak{}INTERVAL\_\allowbreak{}MS}) & 60000 & Status publication interval. \\ \hline
\texttt{arrowhead.\allowbreak{}ssl.\allowbreak{}crl-\allowbreak{}file} & (unset) & CRL checked for every client certificate, reloaded on change, fail closed (standalone deployment: \texttt{/\allowbreak{}crl/\allowbreak{}ca.\allowbreak{}crl}). \\ \hline
\end{longtable}

Secrets (keystore and identity passwords) must be supplied by the deployment; the bundled \texttt{application.\allowbreak{}yml} still has development fallbacks, which the standalone configuration overrides: it reads them from per-deployment Compose file secrets under \texttt{/\allowbreak{}run/\allowbreak{}secrets}.

\subsubsection{Data stored by the system}

Models and observations in memory.

\subsubsection{Non functional properties}

Small in-memory models; not for production learning.

\subsubsection{Stateful or stateless}

Volatile.

\subsection{Important Delimitations}
\label{sec:delimitations}

Experimental; no persistence.

\newpage

\section{Services}
\label{sec:services}

This section describes the produced and consumed services. Each produced service is specified by an SD document and one IDD document per interface.

\subsection{Produced services}

\noindent\begin{longtable}{| p{5.3cm} | p{1.3cm} | p{2.6cm} | p{1.6cm} | p{3.6cm} |} \hline
\rowcolor{gray!33} Service & Version & Interfaces & SD & IDD \\ \hline \endhead
\texttt{aiGym} & 1.0.0 & HTTP(S) & \cite{aig-aiGym-sd} & \cite{aig-aiGym-http-idd} \\ \hline
\texttt{sysMonitor} & 1.0.0 & AMQP(S) & \cite{aig-sysMonitor-sd} & \cite{aig-sysMonitor-amqp-idd} \\ \hline
\end{longtable}

The HTTP interface uses \texttt{generic\_https} with \texttt{CERT\_AUTH}; \texttt{GET /health} is not registered.

\subsection{Consumed services}

\noindent\begin{longtable}{| >{\raggedright\arraybackslash}p{2.9cm} | >{\raggedright\arraybackslash}p{3.1cm} | >{\raggedright\arraybackslash}p{2.5cm} | p{3.6cm} | >{\raggedright\arraybackslash}p{2.9cm} |} \hline
\rowcolor{gray!33} Provider & Service & Operations & Use & Condition \\ \hline \endhead
ServiceRegistry & systemDiscovery, serviceDiscovery & register, lookup & Registration; core lookup. & always \\ \hline
Authentication & identity & identity-login, identity-verify & Own token; caller verification. & always \\ \hline
Consumer\allowbreak{}Authorization & authorization & verify & Per-call authorization. & always \\ \hline
Dynamic\allowbreak{}Service\allowbreak{}Orchestration & service\allowbreak{}Orchestration & pull & sysMonitorData. & always \\ \hline
SysMonitor & sysMonitorData & subscribe & Observation stream. & always \\ \hline
AMQP broker (SysMonitor) & sysMonitor (publisher) & echo, power, status, trainingSample, inference & Own observations. & always \\ \hline
\end{longtable}

Only Authentication's \texttt{identity} and the orchestrator are looked up in the ServiceRegistry; other services are reached through DynamicServiceOrchestration (results cached), and every such call needs a ConsumerAuthorization grant.

\newpage

\section{Security}
\label{sec:security}

The system can run in an insecure mode (\texttt{server.ssl.enabled=false}, plain HTTP, identity declared by the requester) for development, and in the Arrowhead secure mode with HTTPS and X.509 certificates following the Arrowhead certificate profile (\texttt{<System>.<Cloud>.<Organization>.arrowhead.eu}). Certificates not issued under the cloud CA are refused by the TLS handshake. The standalone deployment of this folder (\texttt{config/\allowbreak{}application.\allowbreak{}properties} and \texttt{compose.yaml}) runs the secure mode.

\subsection{Security model}


\begin{itemize}
  \item \textbf{Protocols:} HTTPS (TLS 1.3 and 1.2) with mandatory client certificates; MQTT over TLS; AMQP over TLS for \texttt{sysMonitor}.
  \item \textbf{Data protection:} transport encryption by TLS; the database connection uses TLS with server verification in the standalone deployment. Secrets are kept out of the container environment (Compose file secrets).
  \item \textbf{System authentication:} mutual TLS with a system or operator certificate of the cloud and an identity token of the same system, verified with Authentication; declared identities are not accepted. Client certificates are checked against the CRL.
  \item \textbf{Authorization of produced services:} per operation by caller lists in the implementation and, where stated, by ConsumerAuthorization.
  \item Mutual TLS, identity token bound to the certificate, Consumer\allowbreak{}Authorization per operation.
\end{itemize}

For the Arrowhead certificate profile see github.com/eclipse-arrowhead/documentation.


\bibliographystyle{IEEEtran}
\bibliography{bibliography}

\newpage

\section{Revision History}
\subsection{Amendments}

\noindent\begin{tabularx}{\textwidth}{| p{1cm} | p{3cm} | p{2cm} | X | p{4cm} |} \hline
\rowcolor{gray!33} No. & Date & Version & Subject of Amendments & Author \\ \hline

1 & 2026-10-22 & \arrowversion & Document created from the Arrowhead LaTeX templates for the 5.2.2 Java Spring implementation, including the local extensions of this workspace. & Jerker Delsing \\ \hline
2 & 2026-10-30 & \arrowversion & Checked against the primary Java source tree (Eclipse-Arrowhead-java): standalone deployment, configuration and security statements updated. & Claude Opus 5.5 (draft) \\ \hline
\end{tabularx}

\subsection{Quality Assurance}

\noindent\begin{tabularx}{\textwidth}{| p{1cm} | p{3cm} | p{2cm} | X |} \hline
\rowcolor{gray!33} No. & Date & Version & Approved by \\ \hline

1 & & \arrowversion  &  \\ \hline

\end{tabularx}

\end{document}
